\documentclass[10pt,a4paper]{article}
\usepackage[utf8]{inputenc}
\usepackage[T1]{fontenc}
\usepackage{lmodern}
\usepackage[french]{babel}
\usepackage[top=1.0cm, bottom=1.0cm, left=1.5cm, right=1.5cm]{geometry}
\usepackage{enumitem}
\usepackage{titlesec}
\usepackage{parskip}
\usepackage{xcolor}
\usepackage{hyperref}
\hyphenation{SecNumacadémie authentification PostgreSQL}
\hypersetup{hidelinks}
\pagestyle{empty}
\setlength{\parskip}{2pt}
\definecolor{accentcolor}{RGB}{30, 60, 120}
\titleformat{\section}{\color{accentcolor}\normalfont\bfseries\normalsize}{}{0em}{}[\color{accentcolor}\titlerule]
\titlespacing{\section}{0pt}{4pt}{2pt}
\setlist[itemize]{leftmargin=1.2em, itemsep=0pt, topsep=1pt, parsep=0pt}
\newcommand{\cventry}[4]{%
  \noindent\textbf{#1}\hfill\textit{\small #2}\\
  \noindent\textit{\small #3}%
  \ifx&#4&\else\\\noindent{\small #4}\fi
  \vspace{2pt}
}
\begin{document}

\begin{center}
  {\LARGE\bfseries Bruce TUMPA MADILA}\\[3pt]
  {\color{accentcolor}\normalsize
   Cybersécurité \textbar{} Réseaux \textbar{} SI \textbar{}
   Data \& IA \textbar{} Développement logiciel}\\[2pt]
  \small Contrat d'apprentissage\quad|\quad
         Disponible sept. 2026\quad|\quad
         Île-de-France et France entière\\[2pt]
  \small
  Schiltigheim (67)\quad|\quad
  07~49~19~98~94\quad|\quad
  \href{mailto:bruce.madila@yahoo.fr}{bruce.madila@yahoo.fr}\\[1pt]
  \href{https://linkedin.com/in/bruce-tumpa-madila-a4a184223}{LinkedIn}
  \quad|\quad
  \href{https://github.com/epsi12-lab}{GitHub}
  \quad|\quad
  \href{https://epsi12-lab.github.io/Portofolio}{Portfolio}
\end{center}
\vspace{2pt}

\section{Profil}
Je sécurise des infrastructures réseau, développe des applications distribuées
et analyse des données pour contribuer à vos projets numériques.
Je cherche une alternance en Master pour progresser dans un environnement
exigeant et m'y investir pleinement.

\section{Formation}
\cventry{Licence Informatique}{2022~--~2026}{%
Université de Strasbourg}{%
Algorithmique, systèmes d'exploitation, réseaux, bases de données,
développement logiciel.}

\vspace{3pt}
\cventry{Diplôme de niveau Bac+5~--~Mathématiques,
option Réseaux Informatiques}{2015~--~2020}{%
Université Pédagogique Nationale, Kinshasa (R.D.~Congo)}{%
Mathématiques pures, modélisation, réseaux informatiques.}

\section{Projets techniques}
\cventry{Infrastructure réseau sécurisée~--~Lab GNS3}{2026}{%
VyOS, pfSense, OSPF, VRRP, dnsmasq, Docker}{}
\vspace{-4pt}
\begin{itemize}
  \item Déployé une architecture multi-sites : 4 VLAN, routage OSPF,
    haute disponibilité VRRP et pfSense (refus explicite)~--~délai
    de bascule diagnostiqué par analyse des caches ARP.
\end{itemize}

\cventry{Material Chess~--~Projet intégrateur L3}{2026}{%
Kotlin, Ktor, WebSockets, JWT, PostgreSQL, Redis~--~\href{https://chess.ailurus.fr/}{Démo}}{}
\vspace{-4pt}
\begin{itemize}
  \item Conçu en équipe une architecture microservices (REST API,
    WebSockets, JWT) déployée publiquement~--~sessions sécurisées,
    zéro perte de session constatée.
\end{itemize}

\cventry{Telco Churn Pipeline~--~Lab Data/IA}{2026}{%
Python, Pandas, scikit-learn, matplotlib/seaborn~--~\href{https://github.com/epsi12-lab/telco-churn-pipeline}{GitHub}}{}
\vspace{-4pt}
\begin{itemize}
  \item Pipeline ML de prédiction du churn télécom (IBM) : EDA, feature
    engineering, modélisation supervisée et évaluation F1/ROC-AUC.
\end{itemize}

\cventry{Domaine Active Directory durci~--~Lab GNS3}{2026}{%
Windows Server 2022, AD DS, GPO, ANSSI, PowerShell~--~\href{https://github.com/epsi12-lab/gns3-active-directory}{GitHub}}{}
\vspace{-4pt}
\begin{itemize}
  \item Déployé un domaine AD DS (Win.~Server~2022) avec DNS intégré,
    OU et groupes de sécurité~--~durci par 10 mesures ANSSI (GPO,
    mots de passe, audit Event~4624/4625).
\end{itemize}

\cventry{Likelemba~--~Application mobile d'épargne
collaborative}{En cours}{%
Flutter, Dart, PostgreSQL}{}
\vspace{-4pt}
\begin{itemize}
  \item Développé une application mobile multiplateforme
    (Flutter/Dart) avec persistance PostgreSQL~--~UI fonctionnelle
    Android et iOS depuis un seul codebase.
\end{itemize}

\section{Expérience professionnelle}
\cventry{Opérateur logistique}{CDD\textbar{} 05/2022~--~12/2025}{%
LIDL, direction régionale~--~Strasbourg}{%
Application rigoureuse des procédures qualité~--~zéro incident signalé
sur 3 ans d'activité ; délais maintenus en coordination avec l'équipe.}

\cventry{Stage~--~Systèmes d'information}{Stage\textbar{} 08/2018}{%
Office Congolais de Contrôle, département SI~--~Kinshasa}{}

\section{Compétences techniques et comportementales}
\begin{tabular}{@{}p{3.8cm}p{10.7cm}}
\textbf{Cybersécurité}
  & pfSense, JWT, \mbox{formation SecNumacadémie ANSSI}, triade CIA \\
\textbf{Protocoles réseau}
  & OSPF, EIGRP, VLAN, VRRP, TCP/IP, IPv4/IPv6, DHCP, DNS \\
\textbf{Systèmes}
  & Linux, VyOS, Cisco IOS, Windows Server 2022, AD/GPO \\
\textbf{Développement \& APIs}
  & Kotlin, Dart, Bash, REST API, WebSockets, Ktor, Flutter \\
\textbf{Data \& IA}
  & Python, Pandas, scikit-learn, matplotlib, ML supervisé \\
\textbf{Outils \& logiciels}
  & GNS3, Wireshark, Git/GitHub, Docker, Cisco Packet Tracer \\
\textbf{Bases de données}
  & PostgreSQL, Redis, SQLite, SQL \\
\textbf{Soft skills}
  & Rigueur \textperiodcentered{} Analyse et synthèse
    \textperiodcentered{} Autonomie et initiative
    \textperiodcentered{} Travail en équipe \\
\end{tabular}

\section{Certifications, formations et langues}
\begin{tabular}{@{}p{3.8cm}p{10.7cm}}
\textbf{Langues}
  & Français C1 (courant)\quad Anglais B1 \\
\textbf{Formations}
  & \mbox{Formation SecNumacadémie}~--~ANSSI, 2026 ;
    Réseaux~--~Cisco Networking Academy, 2019 \\
\end{tabular}

\end{document}
